\begin{figure}[!h]
    \centering
    \includegraphics[width=\linewidth]{Images/intro.pdf}
    \caption{Overview of Omni-WorldBench. \textbf{Left:} Omni-WorldSuite defines three levels of interactions, each specified by an initial frame and a prompt. \textbf{Right:} Omni-Metrics comprises an evaluation pipeline that measures interaction effect fidelity, generated video quality, camera-object controllability, and spatiotemporal causal coherence, and then employs an MLLM to adaptively fuse these signals into the final AgenticScore.
    }
\label{fig:intro}
\end{figure}




\section{Introduction}
\label{sec:intro}

%%%% 0301 fxk_V2
The world models aim to characterize the temporal evolution of environmental states under given interaction conditions, providing a foundation for counterfactual reasoning, planning, and decision-making \cite{ha2018world}. Taking advantage of recent advances in video generation, this paradigm has increasingly adopted video synthesis as a core implementation pathway. By leveraging high-quality general-purpose video representations to model world dynamics, video-based world models have been widely applied to autonomous driving, embodied intelligence, and game agents, substantially accelerating progress in these domains.

Unlike rapid progress in world model design, the development of dedicated evaluation benchmarks appears to be somewhat lagging. Existing evaluation methods largely rely on conventional video generation metrics, such as FID and FVD, or adopt general-purpose evaluation benchmarks (e.g., VBench \cite{huang2024vbench}). 
Although these metrics are effective in measuring visual fidelity and text–video alignment \cite{liu2024survey}, 
they do not adequately capture the core capability of world models—the ability to generate consistent and plausible responses under varying interaction conditions.


To comprehensively evaluate the interactive response capabilities of world models, we propose a novel benchmark, \textbf{Omni-WorldBench} (Fig.~\ref{fig:intro}).
At its core, we construct a systematic prompt suite, \textbf{Omni-WorldSuite}, designed to thoroughly assess model performance across diverse interaction levels and scenario types.
Specifically, interaction conditions can produce effects confined to a single object, extend to the local environment, or induce global environmental changes. These progressively increasing interaction scopes impose distinct representational and dynamic modeling requirements on world models. Consequently, the evaluation prompts in Omni-WorldSuite are systematically organized around these three hierarchical interaction levels.
Furthermore, since world modeling is a broad and application-dependent research paradigm, existing studies are often grounded in specific domains such as autonomous driving, embodied robotics, and gaming environments. 
To ensure that Omni-WorldSuite is applicable to both general-purpose video generation models and scenario-specific world models, our evaluation prompts also encompass real-world physical settings as well as representative application domains.

% To complement Omni-WorldSuite, we introduce an effective evaluation metric, \textbf{Omni-Metric}, designed to assess the accuracy of world state representations under specific interaction conditions.
% Specifically, Omni-Metric first focuses on the direct impact of interaction actions on world states. We introduce \textit{Interaction Effect Fidelity} to measure whether an action induces the expected effect on the intervened objects, and \textit{Non-Intervention Consistency} to examine whether unaffected objects undergo unintended changes.
% Beyond assessing the final outcomes of interaction actions, Omni-Metric further incorporates \textit{Spatiotemporal Causal Coherence}, which evaluates whether the evolution of world states follows coherent temporal ordering and plausible spatial causal structure.
% Each evaluation dimension is computed using multiple complementary tools. We further propose a novel agent-based framework that adaptively integrates their outputs conditioned on the prompt semantics to produce the final score.

%方宸v2
% To complement Omni-WorldSuite, we establish a comprehensive and effective evaluation protocol, \textbf{Omni-Metric}, designed to holistically assess the fidelity and consistency of world state representations.
% Distinct from prior works that predominantly focus on static visual fidelity \cite{duan2025worldscore, huang2024vbench++}, Omni-Metric first incorporates \textit{Generated Video Quality}, which extends evaluation to dynamic visual quality. This dimension quantifies temporal artifacts often overlooked in existing benchmarks, including frame jitter, motion smoothness, and dynamic intensity, ensuring the generated sequences maintain high perceptual standards over time.
% Beyond visual quality, we evaluate the stability of the world through \textit{Camera-Object Controllability}. This dimension assess the model's control over camera trajectories and unaffected objects; crucially, we further introduce a novel quantitative analysis of scene transitions, a capability never before addressed in prior works \cite{duan2025worldscore,huang2024vbench++}, to rigorously measure continuity across long-horizon generations.
% The core contribution of our framework lies in interaction consistency, which rigorously examines physical adherence and interactivity. With our proposed \textit{Interaction Effect Fidelity} to measure the expected impact of actions on intervened objects, we expand this scope with multiple new quantitative indicators that verify causal physical laws and spatial logic.
% Finally, reflecting the complexity of these dimensions, each metric is computed using complementary tools, and we propose a novel agent-based framework that adaptively integrates their outputs conditioned on prompt semantics to produce a unified, reliable score.

To complement Omni-WorldSuite, we establish a comprehensive and effective evaluation protocol, \textbf{Omni-Metric}, designed to holistically assess the fidelity and consistency of world state representations. Distinct from prior works that predominantly focus on static visual fidelity \cite{duan2025worldscore, huang2024vbench++}, Omni-Metrics explicitly extends the evaluation toward dynamic, controllable, and interaction-aware generation, which are essential to world models. Specifically, Omni-Metrics evaluates models from three complementary aspects. First, \textit{Generated Video Quality} extends evaluation beyond static appearance to dynamic perceptual quality, measuring temporal flickering, motion smoothness, content alignment, and dynamic degree to capture the visual coherence of generated sequences over time. Second, \textit{Camera-Object Controllability} assesses whether a model can follow explicit camera instructions while maintaining coherent object behavior, and further evaluates long-horizon continuity through a novel scene transition metric, \textit{Transitions Detect}. Third, \textit{Interaction Effect Fidelity} targets the core capability of interactive world modeling by examining whether actions induce the expected effects on intervened objects in a physically plausible and causally consistent manner, supported by quantitative indicators of action-effect correspondence, physical principles, and spatial logic. Since these dimensions are heterogeneous yet complementary, we further introduce an agent-based aggregation framework that adaptively combines outputs from multiple evaluation tools according to prompt semantics, yielding a unified overall metric, \textbf{AgenticScore}, for more reliable evaluation.


Finally, we conduct a systematic evaluation of 18 representative world models, and the results comprehensively reveal the performance boundaries and limitations of current models in interactive response capabilities.
Further human alignment studies demonstrate that Omni-Metric  aligns well with human preferences, validating its effectiveness in assessing world model performance.
Our key contributions are as follows: 
\begin{enumerate}
\item To address the critical absence of standardized evaluation protocols, we introduce \textbf{Omni-WorldBench}. To the best of our knowledge, this is the \textbf{first} benchmark dedicated to assessing the interactive response capabilities of world models, offering a \textbf{comprehensive and holistic} evaluation framework rather than a narrow capability test.
\item We establish a rigorous evaluation infrastructure comprising \textbf{Omni-WorldSuite}, a hierarchical prompt suite spanning diverse interaction levels and scenario types, and \textbf{Omni-Metric}, an agent-based evaluation protocol that quantitatively measures the impact of actions on both final outcomes and intermediate state transitions.
\item We conduct a comprehensive evaluation of 18 video generation models and world models, systematically analyzing their performance. Our findings unveil critical limitations in the 4D interactivity capabilities of current world models, highlighting key areas for improvement. Additionally, we propose a curated benchmark, offering to guide and accelerate future advancements in 4D world model generation.
% \item We conduct a comprehensive evaluation of 18 video generation models and world models, systematically analyzing their performance. Our findings unveil critical limitations in the 4D interactivity capabilities of current world models, highlighting key areas for improvement. Additionally, we propose a curated benchmark and human-alignment evaluation metrics, offering to guide and accelerate future advancements in 4D world model generation.
\end{enumerate}


